\section{Initial regularity: Proof of Theorem \ref{th:reg:1alpha}}\label{sec:initial}
Note that if we
consider the constant factors to be irrelevant with respect to the mathematical argument, for the sake of simplicity
we will omit them in the calculations, writing $\aleq , \ageq , \aeq$ instead of $\leq, \geq$ and $=$.

Most techniques for dealing with critical partial differential equations 
of fractional order have been developed for equations on the whole Euclidean space. For that reason, we
prefer working on the real line over working on the circle.

We will show that for every $u \in \mathbb R / \mathbb Z$ we have
\begin{equation*}
 \lapv g' \in L^{p} ((u-1/20,u+1/20))
\end{equation*}
for a $p>2$. Due to the invariance of the problem under shifting the parametrization, it is enough to show this 
for $u=1/2$, i.e.
\begin{equation} \label{eq:NeedToProveRegOnlyLocally}
 \lapv g' \in L^{p} ((9/20,11/20)).
\end{equation}

Let $\laps{s} \equiv (- \partial_x^2)^{s/2} \equiv (-\lap)^{\frac{s}{2}}$ be the fractional Laplacian on $\R$. The inverse, $\lapms{s} \equiv (-\lap)^{-\frac{s}{2}}$, $s \in (0,1)$, is the Riesz potential by
\begin{equation}\label{eq:fraclap-}
 \lapms{s} f(x) = c_s \intl_{\R} \frac{f(y)}{\abs{x-y}^{1-s}}\ dy.
\end{equation}
In case of positive powers $s$ of the Laplacian $\laps{s}$, $s\in(0,1)$, we use the corresponding formula
\begin{equation}\label{eq:fraclap+}
 \laps{s} f(x) = \tilde c_s \intl_{\R} \frac{f(y)-f(x)}{\abs{x-y}^{1+s}}\ dy.
\end{equation}
For a detailed introduction to the fractional Laplacian we refer to, e.g., \cite{DiNezza}, \cite[Section~2.5]{Schikorra-Doktor}.

To switch from the circle to the real line, we interpret
functions on $\mathbb R / \mathbb Z$ as functions on $\mathbb R$ which are periodic 
with period $1$. We then choose a cutoff function $\eta \in C_0^\infty(\R)$, 
$\eta \equiv 1$ on $[-3,4]$ and consider
\begin{equation}\label{eq:g=g}
  g = \eta \cdot \gamma
\end{equation}
instead of $\gamma$.
We will show that $\lapv g' \in L^p((17/40,23/40),\R^\tdim)$,
for some $p > 2$ to conclude the proof.

Before we begin to outline the structure of the proof, let us shortly recapitulate the notion of 
Lorentz spaces and the main properties we are going to use in this article. For a measurable 
function $f: \Omega \rightarrow \mathbb R$ and $\Omega \subset \mathbb R$
one considers the \emph{decreasing rearrangement}
\begin{equation*}
 f^\ast (t):= \inf\left\{s>0: {\mathcal L}^1(\{|f|> s\})\leq t\right\}
\end{equation*}
where $\mathcal L^1$ denotes the Lebesgue measure. We define
\begin{equation*}
 |f|_{(p,q),\Omega} := \begin{cases}
                      \left(\int_0^\infty (t^{1/p} f^\ast (t)) ^q \frac {dt}{t}\right)^{1/q}  
			&\text{if } p,q \in[1, \infty),\\
                      \sup_{t>0} t^{1/p} f^\ast (t)&\text{if } q=\infty.
                     \end{cases}
\end{equation*}

To prevent technical problems, unless $p \in (1,\infty)$
we will only take the spaces  $L^{1,1} = L^1$ and $L^{\infty,\infty} =  L^\infty$
into consideration.

Furthermore, $\abs f_{(p,p)}\approx\norm f_{L^p}$ for all $p\in[1,\infty]$.
If $\Omega = \mathbb R$ we will omit $\Omega$ in the notation.
Though $|\cdot|_{(p,q), \Omega}$ is not a norm, as it does not obey the triangle inequality, there
is a norm $\|\cdot\|_{(p,q), \Omega}$ on the Lorentz spaces which is equivalent to $|\cdot|_{(p,q),\Omega}$.
These norms satisfy a \emph{H\"older inequality}, i.e., for $p_1,p_2,p \in [1, \infty)$
and $q_1,q_2,q \in [1,\infty]$ with $1/p_1 + 1/p_2 = 1/p$ and $1/q_1 + 1/q_2 = 1/q$ we
have 
\begin{equation}\label{eq:hoelder}
  \|fg\|_{(p,q), \Omega} \aleq \|f\|_{(p_1,q_1), \Omega}\ \|g\|_{(p_2,q_2), \Omega}.
\end{equation}
For $p_1, p_2, p \in (1,\infty)$ and $q_1,q_2 \in [1, \infty]$
with $1/p_1 + 1/p_2 = 1 /p + 1$ and $1/q_1 +1/q_2 = 1/q$ we have the \emph{Young-O'Neil inequality} \cite{Hunt66}
\begin{equation}\label{eq:young}
 \|f \ast g\|_{(p,q), \Omega} \aleq \|f\|_{(p_1,q_1),\Omega}\; \|g\|_{(p_2,q_2), \Omega}.
\end{equation}
Furthermore, we have the \emph{Sobolev inequality} 
\begin{equation*}
 \|\lapms{s} f\|_{(p^\ast,q)} \aleq \|f\|_{(p,q)}
\end{equation*}
for all $s\ge0$, $p\in [1,\infty), q \in [1, \infty]$ and $p^\ast :=  \frac {p}{1-sp} \in [1,\infty)$.
Further information and proofs can be found in~\cite{Hunt66,GrafakosC,Tartar2007}.

The main reason for using Lorentz spaces in the context of critical equations, i.e. equations to which
standard Gagliardo-Nirenberg-Sobolev embeddings cannot be applied to gain regularity, is the following fact.
Although for functions $f$ the $L^2$-norm of $\lapv  f$ does not control the 
$L^\infty$-norm of $f$, the $L^{2,1}$ norm does, i.e. we have the estimate
\begin{equation*}
  \|f\|_{\infty} \aleq \|\lapv   f\|_{(2,1)}.
\end{equation*}
We will also need this in the more general form
\begin{equation*}
 \|f\|_{\infty} \aleq \|\laps{s} f\|_{(\frac 1 s,1)} \quad \mbox{for all $s\in(0,1)$.}
\end{equation*}
% 
% 
In order to prove our regularity result, we will prove a Dirichlet growth theorem for the 
weak $H^{1/2}$-energy of $\gamma'$ on balls in a manner comparable to \cite{DR11Sphere,SchikorranharmSphere}
 -- which are as well in the setting of sphere-valued mappings.
In contrast to these papers, the techniques from \cite{DR11Man,DL11MannD,SchikorraIntBoundFrac11}
deal with a more general setting, but have to work (as we will here) with estimates
of the $L^{2,\infty}$-norm instead of the $L^2$-norm.
Note nevertheless that our right-hand side is very different from their's. In order to obtain the estimates
of the norms $\vrac{\lapv \dsolnt}_{(2,\infty)}$ on small balls, we will have to use new
arguments.

To prove the regularity theorem, we begin with an approach appearing in \cite{DR11Sphere, SchikorranharmSphere} and
divide $\lapv g'$ into the part parallel to $g'$ (and thus \emph{normal} to the sphere $\S^{n-1}$) and
the term perpendicular to $g'$ (\emph{tangential} to the sphere). More precisely, we use that for $p \in [1,\infty)$, $q \in [1,\infty]$
\begin{equation}\label{eq:differencegtanggorth}
 \vrac{\lapv \dsolnt}_{(p,q),B_r} \aleq \vrac{\sbrac{\dsolnt,\lapv \dsolnt}}_{(p,q),B_r} 
  + \sup_{\omega_{ij}} \vrac{\dsolnt_j \omega_{ij}\lapv \dsolnt_i}_{(p,q),B_r},
\end{equation}
where the supremum is over $\omega_{ij} = -\omega_{ji} \in\set{-1,0,1}$.
For a detailed version of this linear algebraic fact, the interested reader is referred to the appendix of \cite{DLSnalphaSphere}.

\subsection{Estimate of the normal part}

We have
\begin{equation}\label{eq:orthogpartHdec}
\sbrac{\dsolnt, \lapv \dsolnt} = -\fracm{2} H_{\frac{1}{2}}(\dsolnt,\dsolnt) + \fracm{2} \lapv \abs{\dsolnt}^{2},
\end{equation}
where
\begin{equation}\label{eq:H}
 H_{s}(a,b) := \laps{s}(ab) -  a\ \laps{s} b -  b\ \laps{s} a.
\end{equation}
Note that for any $s \in (0,1)$, we have
\begin{equation}\label{eq:Dg2}
\laps{s} \sabs{\dsolnt}^{2} \ \refeq{g=g} \ \laps{s} {\eta}^{2} + 2 \laps{s} ({\eta}' \g' g)
+ \laps{s}(|{\eta}'|^2 |\g|^2)
\in L^{\infty}((0,1)). 
\end{equation}
In fact, $\laps{s}\eta^2\in L^\infty$ by interpolation inequalities. For the remaining terms
we use the quasi-locality, Lemma~\ref{lem:QuasiLocality}, and the support of $\eta$ and $\eta'$.

As in \cite{SchikorraIntBoundFrac11} we will use pointwise estimates for $H_s$ and some quantitative version of
the quasi-locality to estimate the normal part of 
$\lapv   g'$:  

\begin{lemma}[Normal part]\label{la:tangentialpart}
For any $s \in [0,\fracm{2})$ there exists $\theta > 0$ such that for any $B_r \subset [0,1]$, $\Lambda > 4$
\begin{equation}\label{eq:orthogpartest}
 \vrac{\laps{s}\sbrac{ \dsolnt, \lapv \dsolnt }}_{(\frac{2}{1+2s},\infty),B_r} 
  \aleq \vrac{\lapv \dsolnt}_{(2,\infty),B_{\Lambda r}}^2 + \ \Lambda^{-\theta} \ \vrac{\lapv \dsolnt}_{(2,\infty),\R}\ 
  \ \sum_{k=1}^\infty 2^{-\theta k}\ \vrac{\lapv \dsolnt}_{(2,\infty),B_{2 ^k\Lambda r}}+r^{\frac12+s}.
\end{equation}
\end{lemma}

For the readers' convenience, a proof will be given in the appendix.

\subsection{Estimate of the tangential part}

It then remains to estimate the part normal to $g'$ (tangential to the sphere), i.e.  for $\omega_{ij} = - \omega_{ji} \in \set{-1,0,1}$, 
$1 \leq i,j \leq n$ we need to estimate 
suitable norms on small balls of the term $\dsolnt_j \omega_{ij} \lapv \dsolnt_i$. 

We have
\begin{align}
\intl_{\R} \dsolnt_j \omega_{ij} \lapv \dsolnt_i\ \lapv \varphi &=
 \intl_{\R} \lapv \dsolnt_i\ \lapv \brac{\omega_{ij} \dsolnt_j \varphi} 
  -\intl_{\R} \lapv \dsolnt_i\ \omega_{ij} \lapv \dsolnt_j\ \varphi 
  -\intl_{\R} \lapv \dsolnt_i\ \omega_{ij} H_{1/2}(\dsolnt_j,\varphi) \nonumber\\
 & = \intl_{\R} \lapv \dsolnt_i\ \lapv \brac{\omega_{ij} \dsolnt_j \varphi} 
  -\intl_{\R} \lapv \dsolnt_i\ \omega_{ij} H_{1/2}(\dsolnt_j,\varphi)\label{eq:antisympde}
\end{align}
where we have used that due to $\omega_{ij} = - \omega_{ji}$ the second term on the right-hand side of the first line
vanishes. 

The second term can be estimated analogously to similar terms in \cite{DR11Man,DL11MannD,SchikorraIntBoundFrac11} using again quasi-locality together with Sobolev embeddings.

\begin{lemma}[Tangential part] \label{lem:CriticalTerm}
There is $\theta >0$, $s_0\in(0,\frac12)$ such that for all 
$\phi \in C_0^\infty(B_r)$, $\Lambda > 16$, $s \in [0,s_0)$  
we have
\begin{multline*}
 \int \lapv   g_i'\omega_{ij} H_{1/2}(g_j',\phi) dx \\\aleq 
  \|\Dsymbol^{-s+1/2}\phi\|_{(\frac 2 {1-2s},1)}
  \left(\|\lapv  g'\|_{(2,\infty),B_{\Lambda r}}^2 + \Lambda^{-\theta} \|\lapv  g'\|_{(2,\infty),\R}
    \sum_{k=1}^\infty 2^{-\theta (k-1)} \|\lapv   g'\|_{(2,\infty), B_{2^k\Lambda r}} \right).
\end{multline*}
\end{lemma}

A proof is provided in the appendix.

It remains to estimate the first term on the right-hand side of Equation~\eqref{eq:antisympde}
for which we will use the Euler-Langrange equation~\eqref{eq:knotPDE}.

Combining this equation with the formula
\begin{multline} \label{eq:FormulaForNorm}
 \int_{\mathbb R } \lapv f_1'\ \lapv f_2' \\
 = c \lim_{\eps\searrow0}\int_{\mathbb R} \int_{\mathbb R\setminus(-\eps,\eps)} 
  \left( \left \langle f_1'(u), f'_2 (u) \right\rangle- 
  \frac{\left\langle f_1(u+w)-f_1(u), f_2(u+w)-f_2(u)\right\rangle }{w^2} \right) 
 \frac {dw}{w^2} du
\end{multline}
due to He \cite[Proposition 2]{He1999} we get the following estimate which 
contains all the information of the Euler-Langrange equation we need to proceed in the proof:

% It 
% In order to do so, we need to use a connection between this first term and 
% $\lim_{\varepsilon \to 0}Q_\varepsilon(\soln,h)$, which is the content of the next lemma. 
% Roughly speaking, it shows that $\lim_{\varepsilon \to 0}Q_\varepsilon(\soln,\cdot)$
% behaves like $\laph \dsolnt$ in the distributional sense -- up to an error of sub-critical terms. 
% This is an adaption of similar results on the torus, cf. \cite[Lemma~2.3]{He00}, and, for a more general setting, \cite[Proposition~1.4]{ReiterOharaKnot11}:
%
\begin{lemma}[Essential estimate of the Euler-Lagrange equation]\label{la:lhs}
  There is a constant $C < \infty$ such that 
  \begin{equation}\label{eq:el-est}
   \int_{\mathbb R} \left\langle \lapv g', \lapv   \phi \right \rangle 
   \leq C \int_{\mathbb R} |\phi(u)| \Gamma(u) du  + C \|\phi\|_{L^2}
  \end{equation}
 for any $\phi \in C^{\infty}_0 ((4/10,6/10), \mathbb R^n)$ with 
 $\left\langle\phi,\gamma'\right\rangle \equiv 0$ where
 \begin{equation*}
  \Gamma(u):= \int_{(-1,1)^3} \int_{-1/4}^{1/4} \frac{|g'(u)-g'(u+s_2 w)| \ 
  |g'(u+s_3 w)-g'(u+s_4 w)|^2} {|w|^2}\ \ dw\ ds.
 \end{equation*}
\end{lemma}


The heart of the proof of Theorem~\ref{thm:FirstStep} is the 
following pointwise estimate of the most problematic term $\Gamma(u)$, which permits to localize it and which afterwards will be transformed 
into a bound of its $L^1$-norm.



\begin{lemma}[Estimate of the critical term]\label{la:T2est}
We have
\begin{equation}\label{eq:Gammauest}
 \Gamma(u) \aleq \abs{\Dsymbol^{-\frac{1}{12}} {\vert \lapv \dsolnt\vert}(u)}^2\ \Dsymbol^{-\frac{1}{3}}{\vert \lapv \dsolnt\vert}(u)+\Dsymbol^{-\frac{1}{12}} {\vert \lapv \dsolnt\vert}(u)\ \Dsymbol^{-\frac{1}{4}}\abs{\Dsymbol^{-\frac{1}{12}} {\vert \lapv \dsolnt\vert}(u)}^2
\end{equation}
almost everywhere.
\end{lemma}

In fact, Lemmata~\ref{la:lhs} and~\ref{la:T2est} provide the essential new estimates in this article
% which technically have not been implicitely used before
on which our entire reasoning relies crucially.
Using these and the above-mentioned improved Sobolev embeddings for
Lorentz spaces, we are ready to show the following estimate. This will allow us to 
prove a Dirichlet growth of the $L^{2, \infty}$-norm of $\lapv   g'$.

\begin{lemma}\label{la:rhsests}
There exist $R > 0$, $s_0 \in (0,\fracm{2})$, and $\sigma>0$
such that for any $\Lambda > 2$, $s \in [0,s_0)$, 
$B_{\Lambda r} \subset (\frac{5}{10},\frac{6}{10})$, 
$r \in (0,\Lambda^{-1}R)$ and $\varphi \in C_0^\infty(B_r)$ we have
\begin{align*}
 &\int \dsolnt_i \omega_{ij} \lapv \dsolnt_j\ \lapv \varphi \\
 &\aleq \vrac{\lapms{s} \lapv \varphi}_{(\frac{2}{1-2s},1)} \brac{r^\sigma + \vrac{\lapv \dsolnt}_{(2,\infty),B_{\Lambda r}}^2+ \Lambda^{-\theta} \vrac{\lapv \dsoln}_{(2,\infty),\R}\sum_{k=1}^\infty 2^{-\theta k}\ \vrac{\lapv \dsoln}_{(2,\infty),B_{2^k\Lambda r}}}.  
\end{align*}
\end{lemma}


\subsection{Conclusion of the proof of Theorem~\ref{thm:FirstStep}}

Combining Lemma~\ref{la:rhsests} with Proposition~\ref{pr:localcontrolelliptic}, we get 
\begin{align}
 &\vrac{\laps{s}\brac{\dsolnt_i \omega_{ij} \lapv \dsolnt_j}}_{(\frac{2}{2s+1},\infty),B_{r}} \nonumber\\
&\aleq 
\brac{\br{\Lambda^{2}r}^\sigma + \vrac{\lapv \dsolnt}_{(2,\infty),B_{\Lambda^{3}r}}^2+ \vrac{\lapv \dsoln}_{(2,\infty),\R} \Lambda^{-\theta} \sum_{k=1}^\infty 2^{-\theta k}\ 
\vrac{\lapv \dsoln}_{(2,\infty),B_{\Lambda^{3} 2^k r}}} \nonumber
\\ &\quad {}+ \Lambda ^{-\theta}r^{-s} \sum_{k=1}^\infty 2^{-\theta k}\ \vrac{\lapv \dsoln}_{(\frac2{1+2s},\infty),B_{\Lambda 2^k r}}\nonumber\\
&\aleq\br{\rule{0ex}{4ex}\cdots} +
\Lambda^{s-\theta} \sum_{k=1}^\infty 2^{(s-\theta) k}\ \vrac{\lapv \dsoln}_{(2,\infty),B_{\Lambda 2^k r}}
\label{eq:nonorthogpartest1}
\end{align}
for some $\theta >0$ uniformly in $\Lambda$.

Let us first use this for $s = 0$, to get in view of \eqref{eq:differencegtanggorth} 
and Lemma \ref{la:tangentialpart} that for all $\varepsilon>0$ we have for sufficiently small
 $r>0$, $B_r \subset \brac{\frac{4}{10},\frac{6}{10}}$  and big enough $\Lambda$
\[
 \vrac{\lapv \dsolnt}_{(2,\infty),B_r} \leq C\br{\br{\Lambda^{2}r}^\sigma + 
  \varepsilon\ \vrac{\lapv \dsolnt}_{(2,\infty),B_{\Lambda^{3} r}}+ (1+\vrac{\lapv \dsoln}_{(2,\infty),\R}) \Lambda^{-\theta} \sum_{k=2}^\infty 2^{-\theta k}\ \vrac{\lapv \dsoln}_{(2,\infty),B_{2^k \Lambda^{3} r}}}
\]
uniformly in $\Lambda$ and $\varepsilon$.

Let us fix such an $r_0$, and consider the above equation for $\Lambda = 2^{m/3}$,
$r = 2^{-m-k} r_0$ (w.l.o.g. $r_0 = 1$).  Setting $b_0 := \|\lapv   g'\|_{(2,\infty)}$ 
and $b_k:= \|\lapv  g'\|_{(2,\infty),B_{2^{-k}}}$ the above estimate gives
\begin{equation*}
 b_{k+m} \leq C 2^{-\sigma (k+m/3)} + \varepsilon b_{k} + 
 C 2^{-\theta m/3}\br{\sum_{l=1}^{k} 2^{-\theta l}b_{k-l}
 +C2^{-\theta k}}
\end{equation*}
for every $k \in \mathbb N_0$ where $C< \infty$ does not depend on $k$.

Using the iteration argument, Lemma~\ref{lem:IterationLemma}, leads to
\begin{equation} \label{eq:morreynormlapvdsolnt}
 \vrac{\lapv \dsolnt}_{(2,\infty),B_r} \leq C_{\soln}\ r^{\tilde{\sigma}}
\end{equation}
for all $B_{r} \subset \brac{\frac{9}{20},\frac{11}{20}}$
 and $r$ small enough.

% In fact, \eqref{eq:morreynormlapvdsolnt} holds for all $B_{\Lambda r}(x_0)$ with $x_0\in(\frac4{10},\frac6{10})$
% and $r>0$ with a constant $C$ depending on $\dist\br{x_0,\set{\frac4{10},\frac6{10}}}$.
% To see this, consider a given compact set $K\subset\brac{\frac{4}{10},\frac{6}{10}}$ and $x_0\in K$.
% If $B_{\Lambda r}(x_0)\subset\brac{\frac{4}{10},\frac{6}{10}}$ then~\eqref{eq:morreynormlapvdsolnt}
% holds for $r<R$ by the iteration argument. If $r\ge R$ we also obtain~\eqref{eq:morreynormlapvdsolnt}
% where $C$ additionally depends on $R$. If, however, $B_{\Lambda r}(x_0)\not\subset\brac{\frac{4}{10},\frac{6}{10}}$
% then $r>R_0(K)$ for some $R_0(K)$ which again leads to~\eqref{eq:morreynormlapvdsolnt}.

This immediately implies H\"older regularity of the solution $g'$. Instead, let us iterate
the above argument to get $\lapv g' \in L^p_{loc}((9/20,11/20))$, for some $p > 2$.

Let us assume without loss of generality that $\tilde{\sigma} < \frac{\theta}{2}$. 
Then, plugging \eqref{eq:morreynormlapvdsolnt} into \eqref{eq:nonorthogpartest1} one obtains
\[
 \vrac{\laps{s}\brac{\dsolnt_i \omega_{ij} \lapv \dsolnt_j}}_{(\frac{2}{2s+1},\infty),B_{r}} 
 \aleq C_{\Lambda,\soln,s,R}\ r^{\min\{\sigma,\tilde{\sigma}, \tilde{\sigma}-s,\theta/2\}}
\]
for small enough $0<s<(0,\tilde\sigma)$, $r>0$, and $B_r \subset [9/20, 11/20]$.
Here we have used that
$2^{k}\Lambda^{3}r\ge R_{0}$ for some $R_{0}>0$ leads
to $2^{-\theta k}\le 2^{-\theta k/2}\br{\frac{\Lambda^{3}}{R_{0}}}^{\theta/2}r^{\theta/2}$.
So the series in~\eqref{eq:nonorthogpartest1}
converge and are bounded by some small positive power of $r$.


% Hence, for small enough $s>0$ there
% is a $\delta > 0$ such that for all small balls $B_{\Lambda r} \subset (\frac{5}{10},\frac{6}{10})$ we have
% \[
%  \vrac{\laps{s}\brac{\dsolnt_i \omega_{ij} \lapv \dsolnt_j}}_{(\frac{2}{2s+1},\infty),B_{\Lambda^{-1} r}}  \aleq C_{\Lambda,\soln}\ r^{\delta}.
% \]
%Let us choose $s>0$ so small that moreover
% $\delta:=\min\{\sigma,\tilde{\sigma}, \tilde{\sigma}-s,\theta/2\}>0$.

On the other hand, \eqref{eq:morreynormlapvdsolnt} and \eqref{eq:orthogpartest} 
together imply for any small enough $r >0$, $\Lambda := r^{-\fracm{2}}$,
and $B_{r}\subset\brac{\frac{9}{20},\frac{11}{20}}$  
\[
\begin{ma} 
 \vrac{\laps{\fracm{4}}\sbrac{ \dsolnt, \lapv \dsolnt }}_{(\frac{4}{3},\infty),B_r}  
&\overset{\eqref{eq:orthogpartest}}{\aleq}& \vrac{\lapv \dsolnt}_{(2,\infty),B_{r^{\fracm{2}}}}^2 + r^{\frac{\theta}{2}}\ \vrac{\lapv \dsolnt}_{(2,\infty),\R}^2 + r^{\frac34}\\
&\overset{\eqref{eq:morreynormlapvdsolnt}}{\leq}& C_{\soln}\ \brac{r^{\tilde{\sigma}} + r^{\frac{\theta}{2}} + r^{\frac34}}.
\end{ma}
\]
That is, $\laps{s}\brac{\dsolnt_i \omega_{ij} \lapv \dsolnt_j}$, and $\laps{\fracm{4}}\sbrac{ \dsolnt, \lapv \dsolnt }$ both belong locally to a Morrey space
$\mathcal L_{loc}^{(p,q),\lambda}$ on $\brac{\frac{9}{20},\frac{11}{20}}$.  
More precisely, for some $\lambda \in (0,1)$ we have
\[
 \begin{ma}
\laps{s}\brac{\dsolnt_i \omega_{ij} \lapv \dsolnt_j}    &\in&  \mathcal{L}^{(\frac{2}{2s+1},\infty),\lambda}_{loc}\brac{\frac{9}{20},\frac{11}{20}}, \\
\laps{\fracm{4}}\sbrac{ \dsolnt, \lapv \dsolnt } &\in& \mathcal{L}^{(\frac{4}{3},\infty),\lambda}_{loc}\brac{\frac{9}{20},\frac{11}{20}}.
 \end{ma}
\]
% Hence,
% \[
%  \begin{ma}
% \laps{s}\brac{\dsolnt_i \omega_{ij} \lapv \dsolnt_j}    &\in&  \mathcal{L}^{\frac{2}{2s+1}-\varepsilon,\lambda}_{loc}(\frac{9}{20},\frac{11}{20}), \\
% \laps{\fracm{4}}\sbrac{ \dsolnt, \lapv \dsolnt } &\in& \mathcal{L}^{\frac{4}{3}-\varepsilon,\lambda}_{loc}(\frac{9}{20},\frac{11}{20})
%  \end{ma}
% \]
% for all $\varepsilon >0$.

The boundedness of Riesz potentials on Morrey spaces, as shown in \cite{Adams75}, 
implies that for some $p > 2$ 
\[
 \dsolnt_i \omega_{ij} \lapv \dsolnt_j,\ \sbrac{ \dsolnt, \lapv \dsolnt } \in L^{p}_{loc}\brac{\frac{9}{20},\frac{11}{20}}.
\]

Together, using \eqref{eq:differencegtanggorth}, we have shown
\[
 \lapv \dsolnt_j \in L^p_{loc}\brac{\frac{9}{20},\frac{11}{20}},
\]
which finishes the proof of Theorem~\ref{th:reg:1alpha}.\qed